\documentclass[10pt,a4paper]{amsart}
\usepackage[margin=19mm]{geometry}
\usepackage{amsmath,amssymb,mathtools}
\usepackage[scr=rsfs,cal=euler]{mathalfa}
\usepackage{xfrac}
\usepackage[unicode,hidelinks,bookmarks=false]{hyperref}
\newcommand{\ie}{i.e.\ }
\newcommand{\cf}{cf.\ }
\newcommand{\resp}{respectively\ }
\newcommand{\Subsection}{\subsection}
\providecommand{\nobreakdash}{\nobreak\hbox{-}}
\setlength{\emergencystretch}{2em}
\multlinegap0pt
\usepackage[shortlabels]{enumitem}
\usepackage{amssymb}
\usepackage{xcolor}
\usepackage{mathtools}
\newtheorem{lemma}{Lemma}
\newtheorem{proposition}{Proposition}
\DeclareMathOperator{\Br}{Br}
\renewcommand{\P}{{\mathbb P}}
\newcommand{\im}{{\rm Im}}
\title{Pullback Method with Applications to Severi--Brauer Fibrations}
\author{Mridul Biswas, Divyasree C Ramachandran, Biswanath Samanta}
\date{Source: arXiv:2410.15125v4 (CC BY 4.0)}
\begin{document}
\maketitle

\noindent\emph{Source and licence.} Pullback Method with Applications to Severi--Brauer Fibrations. Version arXiv:2410.15125v4 (CC BY 4.0), distributed by arXiv under the Creative Commons Attribution 4.0 International licence, http://creativecommons.org/licenses/by/4.0/, as recorded in the arXiv licence metadata for this exact version and shown on the abstract page https://arxiv.org/abs/2410.15125v4. Full licence notice: \texttt{licenses/arxiv-2410.15125-CC-BY-4.0.txt}.\par\smallskip

\noindent\textit{Editorial scope.} Counted unit: the article's proposition lemma (source0.tex lines 505--516). The article's complete original proof of it is delivered with it (source0.tex lines 519--526, one \texttt{proof} environment of 8 lines). No mathematics is written, repaired or re-derived: the statement and the proof are the article's own lines, in the article's own notation, and no retained line is substituted or re-flowed.  Retained uncounted as the source-local context the proof needs: lines 269--277; lines 290--296; lines 492--494.  Nothing was omitted from any retained span, so the package carries no omission record.  No retained line contains a citation, so the wrapper carries no bibliography and no bibliography entry.  No retained line contains a figure environment, an image-inclusion command or drawing code, so no image is delivered and the package declares no asset dependency.  Editorial formatting only, in addition: an AMS \texttt{amsart} preamble that restates the article's own declarations and macros used by the retained lines, namely the environments \texttt{lemma}, \texttt{proposition} on separate counters, together with the standard packages those lines need.  The retained environments print in this document as Lemma 1, Proposition 1 in order of appearance, whereas the article prints the same items with the numbering of its own full document.  The complete native source of the article as distributed in the arXiv e-print for this exact version is bundled unchanged as \texttt{sources/167105/source0.tex}.
	\begin{lemma}\label{SB variety and rational point}
		Let $A$ be a central simple algebra of degree $n$ over a field $k$ and $X$ be the associated Severi--Brauer variety over $k$. The following are equivalent.
		\begin{enumerate}[label=\arabic*.]
			\item $X$ has $k$-rational point.
			\item $X$ is isomorphic to the projective space $\P^{n-1}_{k}$ over $k$.
			\item $A$ splits over $k$. 
		\end{enumerate}
		
	\end{lemma}

    \begin{lemma}\label{split algebra} Let $k$ be a field containing a primitive $m$th root of unity $\zeta_m$ and let $a,b \in k^\times$. Then the following statements are equivalent.
    \begin{enumerate}[label=\arabic*.]
		\item The cyclic algebra $(a,b)_{\zeta_m}$ is split over $k$.
		\item The element $b$ is a norm from the field extension $k(\sqrt[m]{a})$.
		\item The element $a$ is a norm from the field extension $k(\sqrt[m]{b})$.
    \end{enumerate}
    \end{lemma}

\begin{lemma}\label{required polynomial}
    For any given irreducible polynomial $g(x)\in k[x]$ of degree greater than 3, there exists a constant $c\in k$, a non-archimedean place $v$ of $k$, and a polynomial $h(x)\in k[x]$ such that the $v$-adic valuation of $g(c)$ is zero, $g(c)$ is not a $p$-th power in $k_v$, but $f(c)=g(c)h(c)$ is a $p$-th power in $k_v$, where $f(x)=g(x)h(x)$.
\end{lemma}

	\begin{proposition}\label{Main lemma}
		
		
		The variety $\tilde{\pi} : \tilde{X} \to \mathbb{P}_k^1$ constructed above satisfies the following.
		
		\begin{enumerate}[label=\arabic*.]
			\item For all Archimedean places $w$, the Severi--Brauer fibration $\tilde{X}_{k_w}$ is birational to a $\mathbb{P}^{p-1}$ bundle.
			\item There exists a non-Archimedean place $v$ of $k$ not lying over $p$ and a subgroup $H$ of $\Br\tilde{X}$ such that the map $\tilde{\phi}_{v} : \tilde{X}(k_{v}) \to \mathrm{Hom}(H, \mathbb{Q}/\mathbb{Z})$ is nonzero.
			\item For all places $\omega$, we have $0 \in \im (\tilde{\phi}_{\omega})$.	
		\end{enumerate}
		
	\end{proposition}

	\begin{proof}
		(1) The generic fiber of $\tilde{X}$ is the Severi--Brauer variety associated with the cyclic algebra $(a, f)$. After base change to an Archimedean place $w$, $a$ becomes a $p$th power, that is the algebra $(a, f)$ splits over $k_w$. Using Lemma \ref{SB variety and rational point}, we can conclude that the resulting scheme $\tilde{X}_{k_w}$ is birational to a $\mathbb{P}^{p-1}_{k_w}$ bundle. 
		
		(2) Let $v$ be the place coming from Lemma \ref{required polynomial}. Then $f(c)$ is a $p$th power in $k_v$, so the cyclic algebra \((a, f(c)) \) splits in \(\Br k_{v}\), thus the fiber \(\tilde{X}_c\) has a \(k_{v}\)-point, say \(P\). Let \( H\) be the subgroup generated by \((a,f_1)\) inside \(\Br\tilde{X} \). Consider the map $\tilde{\phi}_{v} : \tilde{X}(k_{v}) \to \mathrm{Hom}(H, \mathbb{Q}/\mathbb{Z})$ where \(\tilde{\phi}_{v}(P)\)  sends \((a, f_1)\) to \(\mathrm{inv}_{v}(a, f_1(c))\), as $P$ is above $[c:1]\in \P^1(k_v)$ and $P^{*}(a,f_1)=(a,f_1(c))$. Now, as $f_1(c)$ is not a $p$th power and the \(v\)-adic valuation of \(f_1(c)\) is non-zero, \( k_v(\sqrt[p]{f_1(c)}) \) is a degree-\(p\) unramified extension of \( k_v \), say $L$. Then the valuation of the norm group is given by \( v(N(L^{\times})) = p\mathbb{Z} \). By our choice of \( a \) we have \( v(a) = 1 \), which implies that \( a \) does not belong to the norm group. Consequently, by Lemma \ref{split algebra}, \( (a, f_1(c)) \) does not split in $k_{v}$, thus \(\mathrm{inv}_v(a, f_1(c))\) is non-zero.

		
		(3) Since \( f \) is monic, the variety \( \tilde{X} \) has a \( k \)-point \( Q \) over infinity which implies \( Q^*(a, f_1) = (a, 1) = 0 \). It follows that \( \tilde{\phi}_w(Q) = 0 \) for all places $\omega$.
	\end{proof}

\end{document}
